\section{Conclusion}

This paper evaluated deep learning architectures for RSSI-based passenger movement classification in public transport. Building on classical ML baselines at MCC 0.756, we tested 11+ architectures spanning recurrent, convolutional, transformer, state-space, and mixture-of-experts models. The evaluation used four dataset variants, class-conditioned data augmentation, and exhaustive hyperparameter optimization.

The two-stage DeepStackEnsemble achieved MCC 0.859 and accuracy 89.4\% on combined data, a 13.6\% relative improvement. State-space models reached up to MCC 0.850 with $O(L)$ complexity, confirming their viability as an alternative to transformers for resource-constrained deployment.

Cross-route interference analysis revealed a 9.4\% absolute MCC degradation under multi-device noise. Heatmap analysis identified AA as the dominant interferer and AB as the most vulnerable route. Trajectory visualization showed that concurrent near-AP transmissions erase AB's characteristic 20~dBm RSSI descent. BB proved most noise-resilient because its classification signal comes from temporal flatness, which survives interference-level shifts intact.

Feature importance analysis confirmed and extended prior findings. Classical ML identified early timesteps as most informative. DL additionally captured mid-trajectory motion dynamics through the velocity channel, contributing directly to the 13.6\% improvement. Window deviation emerged as the most informative channel, explaining BB's noise resilience through shape-based classification.

Future work should explore attention mechanisms that learn to down-weight interfered timesteps. Multi-AP fusion could disambiguate near-AP states through spatial diversity. Integration of movement classification with origin-destination matrix inference would enable end-to-end passenger flow modeling. Context-aware architectures that jointly model signal and environmental metadata offer a promising direction for real-world deployment.
